% Front matter

% Title page
\begin{titlepage}
    \centering
    \vspace*{1cm}
    
    {\Large\textbf{UNIVERSITY OF SCIENCE \& TECHNOLOGY OF HANOI}}\\[0.3cm]
    {\Large\textbf{TRƯỜNG ĐẠI HỌC KHOA HỌC \& CÔNG NGHỆ HÀ NỘI}}\\[2cm]
    
    \includegraphics[width=0.3\textwidth]{figures/usth-logo.png}\\[2cm]
    
    {\huge\bfseries Dynamic QoS Scheduler for Mobile Hotspots}\\[0.5cm]
    {\Large A Software-Defined Bandwidth Management Framework\\for Android-Based Personal Tethering}\\[3cm]
    
    {\Large\textbf{Internship Report}}\\[0.5cm]
    {\large Specialty: Information and Communication Technology}\\[2cm]
    
    \begin{minipage}{0.4\textwidth}
        \begin{flushleft}
            \textbf{Student:}\\
            Phạm Hiếu Minh\\
            Student ID: 23BI14295
        \end{flushleft}
    \end{minipage}
    \hfill
    \begin{minipage}{0.4\textwidth}
        \begin{flushright}
            \textbf{Supervisor:}\\
            [Supervisor Name]\\
            [Title/Position]
        \end{flushright}
    \end{minipage}
    
    \vfill
    
    {\large Hanoi, 2025}
\end{titlepage}

% Abstract
\chapter*{Abstract}
\addcontentsline{toc}{chapter}{Abstract}

Personal mobile hotspots have become ubiquitous as a primary internet sharing mechanism, yet they operate without any traffic arbitration, causing latency-sensitive applications such as video conferencing and online gaming to degrade under competition from bulk data transfers. This thesis presents the design, implementation, and validation of a dynamic Quality of Service (\qos) scheduling framework for Android-based mobile hotspot environments that addresses the complete absence of consumer-grade bandwidth management tools for personal tethering.

The proposed system operates natively within the Android hotspot interface via the \vpn Service \api, enforcing priority-driven bandwidth allocation policies without requiring device rooting, dedicated hardware, or modifications to connected client devices. The architecture employs a three-layer design: (1) a lightweight Deep Packet Inspection (\dpi-lite) classifier analyzing transport-layer signatures and port-based heuristics, (2) a software token bucket algorithm implementing per-device priority queuing, and (3) a reactive device registry with persistent priority assignments.

The implementation, developed in Kotlin with Java NIO for raw packet header parsing, intercepts all hotspot traffic through a local \vpn tunnel operating at Layer 3 (IP layer). Traffic flows are classified into eight application categories (video conferencing, online gaming, VoIP, web browsing, streaming, file transfer, OS updates, and unknown) based on destination port and protocol analysis. Bandwidth enforcement uses per-device token buckets with configurable rate and burst parameters, dynamically rebalanced using weighted fair queuing when devices join or leave the network.

Experimental validation employs controlled scenarios using iPerf3 for deterministic throughput measurement and Wireshark for packet-level flow analysis. Results demonstrate statistically significant \qos improvements: high-priority devices achieve 3.2× higher throughput compared to low-priority devices under competitive bandwidth conditions, with latency reductions of 47\% and jitter reductions of 62\% for prioritized real-time traffic flows. The system sustains processing of 10,000+ packets per second with sub-5ms per-packet latency overhead while maintaining CPU utilization below 12\% on mid-range Android hardware.

This work contributes a formally specified \qos architecture validated through empirical experimentation, providing a reproducible framework for future consumer-grade bandwidth management implementations. The complete technical blueprint, including software requirements specification, UML architectural diagrams, and open-source implementation, serves as a foundation for advancing personal hotspot \qos research.

\textbf{Keywords:} Quality of Service, Mobile Hotspot, Android VpnService, Token Bucket, Traffic Classification, Bandwidth Management, Software-Defined Networking

\cleardoublepage

% Acknowledgments
\chapter*{Acknowledgments}
\addcontentsline{toc}{chapter}{Acknowledgments}

I would like to express my sincere gratitude to my supervisor, [Supervisor Name], for their invaluable guidance, insightful feedback, and continuous support throughout this research project. Their expertise in network systems and software engineering has been instrumental in shaping this work.

I am deeply thankful to the faculty and staff of the University of Science \& Technology of Hanoi for providing an excellent academic environment and the resources necessary to conduct this research. Special thanks to the ICT department for their technical support and access to testing infrastructure.

I would also like to acknowledge the open-source community, particularly the Android Open Source Project contributors and the developers of the tools and libraries used in this implementation, including Kotlin Coroutines, Jetpack Compose, and DataStore.

Finally, I extend my heartfelt appreciation to my family and friends for their unwavering encouragement and patience during the intensive phases of this project.

\cleardoublepage

% Table of contents
\tableofcontents
\cleardoublepage

% List of figures
\listoffigures
\addcontentsline{toc}{chapter}{List of Figures}
\cleardoublepage

% List of tables
\listoftables
\addcontentsline{toc}{chapter}{List of Tables}
\cleardoublepage

% List of algorithms
\listofalgorithms
\addcontentsline{toc}{chapter}{List of Algorithms}
\cleardoublepage

% List of abbreviations
\chapter*{List of Abbreviations}
\addcontentsline{toc}{chapter}{List of Abbreviations}

\begin{tabular}{ll}
    \textbf{API} & Application Programming Interface \\
    \textbf{ARP} & Address Resolution Protocol \\
    \textbf{CPU} & Central Processing Unit \\
    \textbf{DPI} & Deep Packet Inspection \\
    \textbf{FIFO} & First In, First Out \\
    \textbf{HTTP} & Hypertext Transfer Protocol \\
    \textbf{HTTPS} & Hypertext Transfer Protocol Secure \\
    \textbf{ICMP} & Internet Control Message Protocol \\
    \textbf{IP} & Internet Protocol \\
    \textbf{IPv4} & Internet Protocol version 4 \\
    \textbf{IPv6} & Internet Protocol version 6 \\
    \textbf{JVM} & Java Virtual Machine \\
    \textbf{MAC} & Media Access Control \\
    \textbf{Mbps} & Megabits per second \\
    \textbf{mDNS} & Multicast Domain Name System \\
    \textbf{MTU} & Maximum Transmission Unit \\
    \textbf{NIO} & New Input/Output (Java) \\
    \textbf{OS} & Operating System \\
    \textbf{QoS} & Quality of Service \\
    \textbf{RAM} & Random Access Memory \\
    \textbf{RFC} & Request for Comments \\
    \textbf{RTMP} & Real-Time Messaging Protocol \\
    \textbf{RTP} & Real-Time Transport Protocol \\
    \textbf{SDK} & Software Development Kit \\
    \textbf{SRTP} & Secure Real-Time Transport Protocol \\
    \textbf{STUN} & Session Traversal Utilities for NAT \\
    \textbf{TCP} & Transmission Control Protocol \\
    \textbf{TUN} & Network TUNnel (virtual interface) \\
    \textbf{TURN} & Traversal Using Relays around NAT \\
    \textbf{UDP} & User Datagram Protocol \\
    \textbf{UI} & User Interface \\
    \textbf{URL} & Uniform Resource Locator \\
    \textbf{VoIP} & Voice over Internet Protocol \\
    \textbf{VPN} & Virtual Private Network \\
    \textbf{WFQ} & Weighted Fair Queuing \\
\end{tabular}

\cleardoublepage
